% ===========================================================================
\chapter{Discuss\~ao}\label{ch:discussao}
% ===========================================================================

Este cap\'itulo interpreta os resultados, examina as causas dos crit\'erios parciais,
avalia o que a media{\c{c}}\~ao por MCP efetivamente agrega, compara o trabalho com a
literatura e explicita as limita{\c{c}}\~oes que permanecem.

\section{Interpreta{\c{c}}\~ao dos resultados}

O resultado mais s\'olido do trabalho n\~ao \'e visual, e sim estrutural: zero n-gons e zero
arestas n\~ao variedade em uma cena de 32.632 faces, com os dois corpos principais
integralmente quadr\^angulares. A relev\^ancia disso est\'a em que a propriedade foi obtida
\emph{por constru{\c{c}}\~ao}, e n\~ao por limpeza posterior. Um modelo gerado por
ferramentas interativas tipicamente requer uma etapa de retopologia --- trabalho manual,
dif\'icil de reproduzir e de auditar. Aqui, a propriedade decorre do algoritmo de tampa e da
estrutura de encadeamento, o que significa que ela se preserva automaticamente em qualquer
reconstru{\c{c}}\~ao.

O segundo resultado relevante \'e a coincid\^encia aritm\'etica entre o previsto pela
estrutura do c\'odigo e o medido no relat\'orio: 46 faixas de 16 faces mais duas tampas de 7
quads resultam em 750, exatamente o valor medido. Essa coincid\^encia tem valor
metodol\'ogico: ela mostra que o relat\'orio mede o que o c\'odigo descreve, e n\~ao
um artefato de avalia{\c{c}}\~ao inserido \emph{a posteriori} pela subdivis\~ao.

O terceiro resultado \'e negativo e igualmente informativo: a inspe{\c{c}}\~ao visual
\emph{comparativa} identificou com clareza um desvio que a medi{\c{c}}\~ao num\'erica
n\~ao poderia detectar. Nenhum dos n\'umeros do relat\'orio --- faces, n-gons, variedade,
envelope --- distingue um modelo com linha de car\'ater de um sem ela. A conclus\~ao
metodol\'ogica \'e que os dois tipos de verifica{\c{c}}\~ao s\~ao complementares e
n\~ao substitu\'iveis, conforme argumentado no Cap\'itulo~\ref{ch:relacionados}.

\section{Sobre os crit\'erios parciais}

\subsection{A silhueta e o problema da aresta viva}

O crit\'erio CA-004 ficou parcial porque falta a linha de car\'ater lateral --- a curva de
raio pequeno que percorre a lateral do ve\'iculo. A causa \'e estrutural e vale a pena
explicit\'a-la, porque ela ilustra um limite do m\'etodo de \emph{loft}.

Um \emph{loft} de esta{\c{c}}\~oes transversais interpola suavemente ao longo do eixo
longitudinal. Se uma fei{\c{c}}\~ao exige uma varia{\c{c}}\~ao de raio pequeno, ela s\'o
s\'o sobrevive \`a subdivis\~ao se estiver sustentada por an\'eis de proximidade
(Cap\'itulo~\ref{ch:fundamentacao}). Sem eles, a subdivis\~ao --- que existe
precisamente para suavizar --- faz o que se espera dela e elimina a descontinuidade. O
modelo, portanto, \'e v\'itima do seu pr\'oprio instrumento: a suavidade que garante a
continuidade dos reflexos \'e a mesma que apaga a linha de car\'ater.

A corre{\c{c}}\~ao \'e conhecida, e n\~ao requer mudan{\c{c}}a de m\'etodo: inserir pares de
esta{\c{c}}\~oes pr\'oximas ao longo da linha pretendida, ou marcar a cadeia de arestas
correspondente com peso de dobra. Os an\'eis de reten{\c{c}}\~ao j\'a implementados no nariz e
na cauda s\~ao a prova de conceito da t\'ecnica. O que faltou foi \emph{aplic\'a-la} \`a
lateral, e n\~ao saber como aplic\'a-la.

H\'a, portanto, uma li{\c{c}}\~ao sobre o pr\'oprio m\'etodo de especifica{\c{c}}\~ao. A
especifica{\c{c}}\~ao exigia ``capturar a linha de cintura ascendente'', mas n\~ao
especificava \emph{como verificar} que a linha estava presente. Um crit\'erio de
aceita{\c{c}}\~ao mais preciso --- por exemplo, ``a vista lateral em material neutro apresenta
uma descontinuidade de curvatura vis\'ivel ao longo de $x \in [-1,0; +1,0]$'' --- teria
orientado a verifica{\c{c}}\~ao desde o in\'icio. Crit\'erios precisam ser
\emph{observ\'aveis}, e n\~ao apenas corretos.

\subsection{O tom da pintura}

O crit\'erio CA-007 ficou parcial por um desvio de tom. A an\'alise identificou tr\^es
fatores concorrentes, mas n\~ao foi poss\'ivel quantificar a contribui{\c{c}}\~ao de cada um.

O primeiro \'e a exposi{\c{c}}\~ao global negativa ($-0,18$), adotada para preservar as altas
luzes. Trata-se de uma decis\~ao de compromisso: sem ela, os reflexos especulares no
para-lama estouravam e mascaravam a leitura de continuidade. O efeito colateral \'e o
escurecimento geral.

O segundo \'e a diferen{\c{c}}a de cadeia de imagem: a refer\^encia foi processada
pela c\^amera e pelo \emph{software} fotogr\'afico; o render passou pelo AgX. Comparar tom
entre duas cadeias diferentes sem um alvo de cor comum \'e, em rigor, uma compara{\c{c}}\~ao
inv\'alida.

O terceiro \'e a escolha do pr\'oprio valor de cor base. O azul foi definido como
\texttt{\#0C1A96} convertido para linear, um valor escuro. Aumentar a luminosidade da cor
base \'e a corre{\c{c}}\~ao mais direta, mas s\'o faz sentido depois de eliminar a
contribui{\c{c}}\~ao de exposi{\c{c}}\~ao e cadeia de imagem.

A li{\c{c}}\~ao metodol\'ogica \'e que um crit\'erio de aceita{\c{c}}\~ao sobre cor
\'e, sem refer\^encia colorim\'etrica, um crit\'erio mal posto. Ele se aproxima mais de uma
opini\~ao fundamentada do que de uma verifica{\c{c}}\~ao. Reconhecer isso \'e mais
\'util que tentar ajustar o tom por tentativa e erro at\'e a coincid\^encia visual.

\section{O que a media{\c{c}}\~ao por MCP agrega}

O MCP foi empregado como meio, e n\~ao como objeto. Ainda assim, a an\'alise do seu papel
produz tr\^es conclus\~oes.

\textbf{A serializa{\c{c}}\~ao \'e o ganho principal.} O protocolo n\~ao torna o
Blender mais capaz: todas as opera{\c{c}}\~oes realizadas seriam poss\'iveis por
\emph{script} executado na linha de comando. O que ele agrega \'e que cada
intera{\c{c}}\~ao \'e uma mensagem estruturada com nome de ferramenta e argumentos. Isso
desloca o registro do processo de um hist\'orico de terminal --- que se perde, e que n\~ao
\'e verificado por ningu\'em --- para um conjunto de mensagens inspecion\'aveis.

\textbf{A verifica{\c{c}}\~ao visual por agente \'e limitada.} A inspe{\c{c}}\~ao mediada por
captura de \emph{viewport} \'e estritamente inferior \`a inspe{\c{c}}\~ao interativa: o
avaliador n\~ao pode orbitar, aproximar nem alterar a ilumina{\c{c}}\~ao para testar
hip\'oteses sobre a superf\'icie. A limita{\c{c}}\~ao \'e intr\'inseca \`a natureza
bidimensional da captura \cite{hartley2004}. Neste trabalho, ela foi compensada por
vistas m\'ultiplas predefinidas e por verifica{\c{c}}\~ao num\'erica, mas continua a
existir.

\textbf{O valor est\'a na reprodutibilidade, n\~ao na velocidade.} O tempo de build
(0,6\,s) \'e irrelevante perto do tempo de render (45\,s para nove vistas). N\~ao h\'a,
neste caso, ganho de velocidade apreci\'avel em rela{\c{c}}\~ao a um fluxo manual bem
executado. O ganho \'e que o artefato pode ser reconstru\'ido identicamente, revisado por
diferen{\c{c}}a de c\'odigo e verificado por medi{\c{c}}\~ao autom\'atica. \'E um ganho
de \emph{qualidade de processo}, e n\~ao de produtividade bruta.

Uma consequ\^encia pr\'atica dessa an\'alise \'e que o uso de agentes via MCP faz mais
sentido em geometria \emph{param\'etrica e repet\'ivel} --- rodas, louvres, grelhas,
conjuntos de pe{\c{c}}as com varia{\c{c}}\~ao --- do que em superf\'icies nas quais o
julgamento est\'etico cont\'inuo \'e dominante. A distribui{\c{c}}\~ao de objetos
apresentada na Tabela~\ref{tab:grupos} \'e sugestiva: cinquenta e dois objetos de roda
foram gerados por uma \'unica fun{\c{c}}\~ao parametrizada, enquanto a superf\'icie do
casco --- um \'unico objeto --- exigiu a maior parte do esfor{\c{c}}o de itera{\c{c}}\~ao.

\section{Compara{\c{c}}\~ao com a literatura}

Frente \`a modelagem procedimental cl\'assica \cite{parish2001, muller2006, mccormack2004},
este trabalho n\~ao prop\~oe gram\'atica de forma nova. Sua contribui{\c{c}}\~ao, nessa
frente, \'e mostrar que a parametriza{\c{c}}\~ao por tabela de esta{\c{c}}\~oes --- um
m\'etodo simples e pouco citado --- \'e suficiente para uma morfologia automotiva
plaus\'ivel quando combinada com subdivis\~ao e controle de densidade.

Frente \`a engenharia reversa geom\'etrica \cite{varady1997, otto2001}, a diferen{\c{c}}a
\'e de entrada e de objetivo. A literatura parte de medi{\c{c}}\~oes densas e visa
manufatura; aqui, parte-se de seis fotografias em perspectiva e visa renderiza{\c{c}}\~ao.
Essa diferen{\c{c}}a desloca o problema de fidelidade microm\'etrica para fidelidade de
morfologia e propor{\c{c}}\~ao --- um regime em que a aus\^encia de blueprints
ortogr\'aficos \'e a principal fonte de erro, conforme discutido nas
amea{\c{c}}as \`a validade.

Frente \`a reconstru{\c{c}}\~ao por imagens \cite{snavely2006, mildenhall2020, kerbl2023}, a
diferen{\c{c}}a \'e de natureza do artefato. Uma reconstru{\c{c}}\~ao por
\emph{Structure-from-Motion} ou por campos de radi\^ancia produziria uma representa{\c{c}}\~ao
visualmente mais fiel sob os pontos de vista observados, mas com malha irregular e sem
fluxo de arestas --- inadequada ao requisito de topologia. \'E poss\'ivel que a
combina{\c{c}}\~ao dos dois m\'etodos seja mais interessante que qualquer um isoladamente:
usar reconstru{\c{c}}\~ao densa para \emph{validar} propor{\c{c}}\~oes e o modelo
procedimental para produzir a malha edit\'avel. Essa combina{\c{c}}\~ao n\~ao foi
explorada aqui e \'e sugerida como trabalho futuro.

Frente aos agentes com ferramentas \cite{schick2023, qin2023, yao2023} e ao pr\'oprio
MCP \cite{mcpspec2024}, a contribui{\c{c}}\~ao \'e um relat\'o de caso com
verifica{\c{c}}\~ao. A literatura descreve a capacidade de invocar ferramentas; este
trabalho registra o que aconteceu quando essa capacidade foi aplicada a um problema com
crit\'erio de aceita{\c{c}}\~ao externo --- e registra tamb\'em que o agente n\~ao
substitui o julgamento est\'etico nem a decis\~ao de engenharia sobre o que verificar.

\section{Limita{\c{c}}\~oes}\label{sec:limitacoes-discussao}

Al\'em das amea{\c{c}}as \`a validade j\'a discutidas na Se{\c{c}}\~ao~\ref{sec:ameacas},
quatro limita{\c{c}}\~oes s\~ao insepar\'aveis do resultado obtido.

\textbf{Aus\^encia de avalia{\c{c}}\~ao independente.} A avalia{\c{c}}\~ao visual foi feita
pelo autor do modelo. Sem avalia{\c{c}}\~ao cega e sem m\'ultiplos avaliadores, n\~ao h\'a
como estimar a concord\^ancia do julgamento.

\textbf{Aus\^encia de refer\^encia colorim\'etrica.} Como discutido, o crit\'erio de cor
n\~ao \'e verific\'avel sem um alvo f\'isico e uma cadeia comum de captura.

\textbf{Verifica{\c{c}}\~ao de continuidade por inspe{\c{c}}\~ao dos reflexos.} A
especifica{\c{c}}\~ao original pedia valida{\c{c}}\~ao por mapas de zebra e isofotas. Essa
valida{\c{c}}\~ao foi declarada fora de escopo por decis\~ao do usu\'ario, o que
significa que a continuidade $G^2$ alegada no material e na geometria n\~ao foi medida ---
apenas observada. Uma afirmativa forte como ``continuidade $G^2$'' exigiria an\'alise
dedicada.

\textbf{Estudo de caso \'unico.} H\'a um \'unico ve\'iculo, um \'unico acabamento e uma
\'unica configura{\c{c}}\~ao de \emph{software}. As conclus\~oes gerais s\~ao
hip\'oteses, e n\~ao resultados.

\section{Implica{\c{c}}\~oes}

Para a pr\'atica de modelagem, o trabalho sugere tr\^es diretrizes.

A primeira \'e \emph{impor as restri{\c{c}}\~oes topol\'ogicas no algoritmo}. Verificar n-gons
ao final do processo e corrigi-los \'e uma batalha perdida; construir \'unicamente
estruturas que n\~ao os produzem \'e uma vit\'oria permanente. A fun{\c{c}}\~ao de tampa por
leque dobrado \'e um exemplo pequeno e reutiliz\'avel dessa postura.

A segunda \'e \emph{separar a verifica{\c{c}}\~ao por natureza da grandeza}. Grandezas
estruturais s\~ao medidas; grandezas perceptuais s\~ao inspecionadas. Misturar as duas
leva tanto a medir o que n\~ao importa quanto a aceitar o que n\~ao se pode medir.

A terceira \'e \emph{tratar a automa{\c{c}}\~ao por agente como infraestrutura, e n\~ao
como substituto do julgamento}. O agente executa, mede e reporta com fidelidade not\'avel;
ele n\~ao decide o que \'e bonito, nem que toler\^ancia \'e aceit\'avel, nem que
desvio merece ser perseguido. Neste trabalho, a an\'alise visual foi o instrumento
decisivo para identificar o \'unico desvio que a medi{\c{c}}\~ao n\~ao alcan{\c{c}}ava.

\section{Contribui{\c{c}}\~oes}

Sintetizando, o trabalho contribui com: (i) um gerador procedimental completo e
reproduz\'ivel de um ve\'iculo esportivo, com topologia controlada por constru{\c{c}}\~ao;
(ii) uma t\'ecnica de tampa de \emph{loft} que produz apenas quads, aplic\'avel a
qualquer superf\'icie fechada por an\'eis; (iii) uma demonstra{\c{c}}\~ao de arquitetura
de agentes mediada por MCP para conte\'udo tridimensional, com o registro do que ela
agrega e do que n\~ao agrega; e (iv) um conjunto de evid\^encia visual comparativa
organizado, com resultado registrado por crit\'erio, incluindo os desvios.
